\documentclass[10pt,a4paper]{amsart}
\usepackage{amsmath,amssymb,amsthm,graphicx}
\usepackage{verbatim}
\newenvironment{reference}{\comment}{\endcomment}
%\newenvironment{reference}{}{}
\newenvironment{slogan}{\comment}{\endcomment}
\newenvironment{history}{\comment}{\endcomment}

% For commutative diagrams we use Xy-pic
\usepackage[all]{xy}

% We use 2cell for 2-commutative diagrams.
\xyoption{2cell}
% Theorem environments.
%
\theoremstyle{plain}
\newtheorem{theorem}[subsection]{Theorem}
\newtheorem{proposition}[subsection]{Proposition}
\newtheorem{lemma}[subsection]{Lemma}

\theoremstyle{definition}
\newtheorem{definition}[subsection]{Definition}
\newtheorem{example}[subsection]{Example}
\newtheorem{exercise}[subsection]{Exercise}
\newtheorem{situation}[subsection]{Situation}

\theoremstyle{remark}
\newtheorem{remark}[subsection]{Remark}
\newtheorem{remarks}[subsection]{Remarks}

\numberwithin{equation}{subsection}

% Macros
%
\def\lim{\mathop{\mathrm{lim}}\nolimits}
\def\colim{\mathop{\mathrm{colim}}\nolimits}
\def\Spec{\mathop{\mathrm{Spec}}}
\def\Hom{\mathop{\mathrm{Hom}}\nolimits}
\def\Ext{\mathop{\mathrm{Ext}}\nolimits}
\def\SheafHom{\mathop{\mathcal{H}\!\mathit{om}}\nolimits}
\def\SheafExt{\mathop{\mathcal{E}\!\mathit{xt}}\nolimits}
\def\Sch{\mathit{Sch}}
\def\Mor{\mathop{\mathrm{Mor}}\nolimits}
\def\Ob{\mathop{\mathrm{Ob}}\nolimits}
\def\Sh{\mathop{\mathit{Sh}}\nolimits}
\def\NL{\mathop{N\!L}\nolimits}
\def\CH{\mathop{\mathrm{CH}}\nolimits}
\def\proetale{{pro\text{-}\acute{e}tale}}
\def\etale{{\acute{e}tale}}
\def\QCoh{\mathit{QCoh}}
\def\Ker{\mathop{\mathrm{Ker}}}
\def\Im{\mathop{\mathrm{Im}}}
\def\Coker{\mathop{\mathrm{Coker}}}
\def\Coim{\mathop{\mathrm{Coim}}}

% Boxtimes
%
\DeclareMathSymbol{\boxtimes}{\mathbin}{AMSa}{"02}

%
% Macros for moduli stacks/spaces
%
\def\QCohstack{\mathcal{QC}\!\mathit{oh}}
\def\Cohstack{\mathcal{C}\!\mathit{oh}}
\def\Spacesstack{\mathcal{S}\!\mathit{paces}}
\def\Quotfunctor{\mathrm{Quot}}
\def\Hilbfunctor{\mathrm{Hilb}}
\def\Curvesstack{\mathcal{C}\!\mathit{urves}}
\def\Polarizedstack{\mathcal{P}\!\mathit{olarized}}
\def\Complexesstack{\mathcal{C}\!\mathit{omplexes}}
% \Pic is the operator that assigns to X its picard group, usage \Pic(X)
% \Picardstack_{X/B} denotes the Picard stack of X over B
% \Picardfunctor_{X/B} denotes the Picard functor of X over B
\def\Pic{\mathop{\mathrm{Pic}}\nolimits}
\def\Picardstack{\mathcal{P}\!\mathit{ic}}
\def\Picardfunctor{\mathrm{Pic}}
\def\Deformationcategory{\mathcal{D}\!\mathit{ef}}
\usepackage{cleveref}
\crefname{originalproposition}{original proposition}{original propositions}
\Crefname{originalproposition}{Original proposition}{Original propositions}
\makeatletter
\newlabel{proposition-loc-qcoh-flat-base-change}{{0771}{source}{}{sourcepointer-1}{}}
\newlabel{proposition-loc-qcoh-flat-base-change@cref}{{[originalproposition][0][]0771}{[1][0][]source}{}{}{}}
\makeatother
\crefname{originallemma}{original lemma}{original lemmas}
\Crefname{originallemma}{Original lemma}{Original lemmas}
\makeatletter
\newlabel{lemma-loc-qcoh-fbc-compute-higher-direct-image}{{0GQH}{source}{}{sourcepointer-2}{}}
\newlabel{lemma-loc-qcoh-fbc-compute-higher-direct-image@cref}{{[originallemma][0][]0GQH}{[1][0][]source}{}{}{}}
\makeatother
\crefname{originalreference}{original reference}{original references}
\Crefname{originalreference}{Original reference}{Original references}
\makeatletter
\newlabel{spaces-cohomology-section-higher-direct-image}{{071Y}{source}{}{sourcepointer-3}{}}
\newlabel{spaces-cohomology-section-higher-direct-image@cref}{{[originalreference][0][]071Y}{[1][0][]source}{}{}{}}
\makeatother
\crefname{originalreference}{original reference}{original references}
\Crefname{originalreference}{Original reference}{Original references}
\makeatletter
\newlabel{coherent-lemma-relative-affine-vanishing}{{01XC}{source}{}{sourcepointer-4}{}}
\newlabel{coherent-lemma-relative-affine-vanishing@cref}{{[originalreference][0][]01XC}{[1][0][]source}{}{}{}}
\makeatother
\crefname{originalreference}{original reference}{original references}
\Crefname{originalreference}{Original reference}{Original references}
\makeatletter
\newlabel{morphisms-lemma-affine-equivalence-modules}{{01SB}{source}{}{sourcepointer-5}{}}
\newlabel{morphisms-lemma-affine-equivalence-modules@cref}{{[originalreference][0][]01SB}{[1][0][]source}{}{}{}}
\makeatother
\title{Affine morphisms: vanishing, exactness and direct sums on LQCoh-fbc}
\author{The Stacks Project Authors}
\date{stacks-project:stacks-cohomology.tex:local-snapshot-20261012; GNU FDL 1.2 or later}
\begin{document}
\maketitle
\noindent\textit{Editorial note.} All delivered mathematical text is the author's own, copied verbatim from the pinned Stacks Project chapter file stacks-cohomology.tex; no statement or proof step was rewritten, reordered or removed.\par\medskip
\noindent\textit{Editorial note.} The counted span cites the lemmas recorded in assembly.external\_references with canonical tags as navigation pointers.\par\medskip
\noindent\textit{Editorial note.} Spot-check note: The algebraic-space-to-scheme reduction is delegated to 'the discussion in Cohomology of Spaces, Section spaces-cohomology-section-higher-direct-image' - a section-discussion reference rather than a numbered lemma.\par\medskip
\noindent\textit{Allowed standard prerequisites (editorial).}\begin{itemize}
\item Higher direct image restriction. Cohomology of Algebraic Stacks, Lemma stacks-cohomology-lemma-loc-qcoh-fbc-compute-higher-direct-image. Reduction to algebraic spaces.
\item Relative affine vanishing. Cohomology of Schemes, Lemma coherent-lemma-relative-affine-vanishing. Scheme-case vanishing.
\item Affine equivalence of modules. Morphisms, Lemma morphisms-lemma-affine-equivalence-modules. Direct sums.
\end{itemize}\medskip
\ifdefined\hypertarget\hypertarget{sourcepointer-1}{}\fi
\noindent\textit{Source reference.} Original proposition 0771: sources/stacks-chow/tags.txt, source lines 9186--9186. Cross-chapter standard prerequisite cited by the author inside the counted span; it is used as a pointer and its own proof is not part of this item.\par\medskip
\ifdefined\hypertarget\hypertarget{sourcepointer-2}{}\fi
\noindent\textit{Source reference.} Original lemma 0GQH: sources/stacks-chow/tags.txt, source lines 20862--20862. Cross-chapter standard prerequisite cited by the author inside the counted span; it is used as a pointer and its own proof is not part of this item.\par\medskip
\ifdefined\hypertarget\hypertarget{sourcepointer-3}{}\fi
\noindent\textit{Source reference.} Original reference 071Y: sources/stacks-chow/tags.txt, source lines 9007--9007. Cross-chapter standard prerequisite cited by the author inside the counted span; it is used as a pointer and its own proof is not part of this item.\par\medskip
\ifdefined\hypertarget\hypertarget{sourcepointer-4}{}\fi
\noindent\textit{Source reference.} Original reference 01XC: sources/stacks-chow/tags.txt, source lines 2513--2513. Cross-chapter standard prerequisite cited by the author inside the counted span; it is used as a pointer and its own proof is not part of this item.\par\medskip
\ifdefined\hypertarget\hypertarget{sourcepointer-5}{}\fi
\noindent\textit{Source reference.} Original reference 01SB: sources/stacks-chow/tags.txt, source lines 2332--2332. Cross-chapter standard prerequisite cited by the author inside the counted span; it is used as a pointer and its own proof is not part of this item.\par\medskip
\begin{lemma}
\label{lemma-loc-qcoh-fbc-affine-direct-image}
Let $f : \mathcal{X} \to \mathcal{Y}$ be an affine morphism of algebraic
stacks. The functor $f_* : \textit{LQCoh}^{fbc}(\mathcal{O}_\mathcal{X}) \to
\textit{LQCoh}^{fbc}(\mathcal{O}_\mathcal{Y})$ is exact and commutes
with direct sums. The functors $R^if_*$ for $i > 0$ vanish on
$\textit{LQCoh}^{fbc}(\mathcal{O}_\mathcal{X})$.
\end{lemma}

\begin{proof}
The functors exist by Proposition \ref{proposition-loc-qcoh-flat-base-change}.
By Lemma \ref{lemma-loc-qcoh-fbc-compute-higher-direct-image}
this reduces to the case of an affine morphism of algebraic spaces
taking higher direct images in the setting of quasi-coherent
modules on algebraic spaces. By the discussion in
Cohomology of Spaces, Section
\ref{spaces-cohomology-section-higher-direct-image}
we reduce to the case of an affine morphism of schemes.
For affine morphisms of schemes we have the vanishing of
higher direct images on quasi-coherent modules by
Cohomology of Schemes, Lemma \ref{coherent-lemma-relative-affine-vanishing}.
The vanishing for $R^1f_*$ implies exactness of $f_*$.
Commuting with direct sums follows from
Morphisms, Lemma \ref{morphisms-lemma-affine-equivalence-modules}
for example.
\end{proof}

\end{document}
